\documentclass[11pt,a4paper]{article}
\usepackage[utf8]{inputenc}
\usepackage[T1]{fontenc}
\usepackage{lmodern}
\usepackage{amsmath,amssymb,amsthm}
\usepackage{mathtools}
\usepackage{geometry}
\geometry{margin=2.5cm}
\usepackage{hyperref}
\usepackage{graphicx}
\usepackage{booktabs}
\usepackage{enumitem}
\usepackage{amsfonts}
\usepackage{amsopn}
\usepackage{doi}

\input{glyphtounicode}
\pdfgentounicode=1

\newtheorem{theorem}{Theorem}[section]
\newtheorem{proposition}[theorem]{Proposition}
\newtheorem{lemma}[theorem]{Lemma}
\newtheorem{corollary}[theorem]{Corollary}
\newtheorem{definition}[theorem]{Definition}
\newtheorem{remark}[theorem]{Remark}
\newtheorem{hypothesis}[theorem]{Hypothesis}

\newcommand{\Rn}{R_n}
\newcommand{\Sn}{\bar{\Sigma}_n}
\newcommand{\dexact}{\hat{\delta}_{\mathrm{exact}}}
\newcommand{\deff}{\delta_{\mathrm{eff}}}
\newcommand{\Reff}{R_n^{\mathrm{eff}}}
\newcommand{\bstar}{\beta^{*}}

\title{Observable-Class Derivation Mismatch and the Heisenberg Transfer No-Go:\\
From Proxy Capacity to Block-Level Span Dynamics}
\author{J\'er\^ome Beau\\
\small Independent Researcher, France; jerome.beau@cosmochrony.org}
\date{2026\\Version 1.15.2}

\begin{document}
\maketitle

\begin{abstract}
The O3--O7 derivation chain proposes, on a changing-degree LPS family, the relation
$\bstar = 1/(\delta + 1/2)$ linking the capacity decay exponent~$\delta$
to the cascade growth exponent~$\bstar$ governing charged-lepton mass ratios.
Papers O12--O14 demonstrated that the exact Weil-block capacity exponent
$\dexact(q) \in [3.6, 4.8]$ is incompatible with the target window
$\delta \in [7.4, 10.6]$, and that this tension is structural rather than
a finite-size artefact.
The present paper identifies the precise origin of the mismatch and proves a
non-transferability result:
the O6 growth equation $\mathrm{d}p/\mathrm{d}n \lesssim \Rn \, c_{\mathrm{BI}} \sqrt{p(n)}$
derives $\bstar$ from a \emph{scalar global} redundancy functional~$\Rn$,
whereas the exact observable measured in O12--O13 is a \emph{block-averaged}
capacity over $(q-1)$ Weil representations of internal dimension~$q$.
We prove a non-equivalence theorem establishing that the block-averaged
exponent $\dexact$ does not, in general, coincide with the scalar exponent~$\alpha$
entering the growth equation.
We define an effective block-weighted redundancy $\Reff$ and show that its exponent cannot
repair the gap by aggregation alone.
The native identity is instead
$\Delta r_n=|S_n|\Sigma_n\asymp \Sigma_n n^{D-1}$.
Neither the Born--Infeld factor $c_{\mathrm{BI}}\sqrt{p}$ nor the changing LPS valence has a
native carrier on the fixed-degree Heisenberg cascade.
Moreover, a constant factor $q^{-\eta}$ cannot alter an intra-$q$ slope, and direct exponent
matching in the transplanted equation does not yield the legacy reciprocal formula.
The former C1/C2 outputs are therefore retained only as phenomenological diagnostics.
Block-level redundancy measurement remains useful for the aggregation question, not as a test
of a native capacity-to-rate transfer.
\end{abstract}

\smallskip\noindent
\textbf{Keywords.}
Observable-class mismatch, growth equation, block-averaged capacity,
scalar redundancy, Weil representations, Heisenberg graphs,
cascade exponent, spectral admissibility.

  \clearpage
\tableofcontents

\clearpage
\subsection*{Notation glossary}

  \noindent
  The four exponents appearing in this paper form a strict hierarchy:
  \begin{center}
    \begin{tabular}{clp{7.5cm}}
      \toprule
      Symbol & Name & Definition \\
      \midrule
      $\dexact$ & measured exponent &
      Power-law fit of the unweighted block average
      $\Sn = \mathbb{E}_c[\Sigma_n^{(c)}] \sim A\,n^{-\dexact}$
      (O12/O13 observable). \\[4pt]
      $\alpha_{\mathrm{dyn}}$ & aggregation exponent &
      Power-law fit of the weighted observable
      $\Reff \sim B\,n^{-\alpha_{\mathrm{dyn}}}$
      (Definition~\ref{def:alphaeff}).
      Satisfies $\alpha_{\mathrm{dyn}} \leq \dexact$
      (Corollary~\ref{cor:nogo}). \\[4pt]
      $\deff$ & legacy endpoint diagnostic (C1) &
      Historical redefinition combining shell growth and an endpoint
      $q$-normalisation term:
      $\deff = \dexact - (D-1) + \log q / \log n^{*}$. \\[4pt]
      $\sigma(q)$ & legacy diagnostic term (C2) &
      Additional $q$-dependent denominator term chosen to quantify the
      displacement required to reach $\bstar \in (0.09, 0.13)$; no native carrier is derived. \\
      \bottomrule
    \end{tabular}
  \end{center}

\section{Introduction}
\label{sec:intro}

\subsection{Position in the O-series}
\label{sec:position}

The admissibility sub-programme has traced a sequence from obstruction
identification (O8--O9~\cite{Beau2026a12,Beau2026a13}) through exact
measurement (O10--O13~\cite{Beau2026a14,Beau2026a15,Beau2026a16,Beau2026a17})
to observable and estimator separation (O14~\cite{Beau2026a18}).
O14 established that the tension between the exact Weil-block capacity
exponent $\dexact(q) \in [3.6, 4.8]$ and the phenomenological target
$\delta \in [7.4, 10.6]$ is not a numerical artefact (refuting hypothesis~S1
of O13) but a structural mismatch (Source~S2).
The conditional endpoint diagnostic retained in O14,
\begin{equation}
\label{eq:O14corrected}
\beta_{\mathrm{diag}} = \frac{1}{\deff + \tfrac{1}{2}} + \epsilon(q),
\end{equation}
uses $\deff = \dexact - \delta_{\gamma,\mathrm{corr}}
- \eta \cdot \log q / \log n^{*}$ only as an endpoint/inter-$q$ estimator.
The legacy diagnostic selects Scenario~S2-B for all tested primes.
At fixed $q$, the factor $q^{-\eta}$ changes the regression intercept, not its slope.

O14 thus separated the measurement question from the
\emph{derivation} question:
does the functional form $\bstar = 1/(\delta + 1/2)$ itself remain valid
when the input observable changes from a scalar global capacity to a
block-averaged Weil capacity?
The present paper answers this question negatively for the corpus-defined objects and isolates
the additional hypotheses a future model would require~\cite{Beau2026sgn}.

\subsection{The central question}
\label{sec:central}

The relation $\bstar = 1/(\delta + 1/2)$ was derived conditionally in
O6~\cite{Beau2026a10} (Proposition~4.2) from the LPS growth equation
\begin{equation}
\label{eq:growth}
\frac{\mathrm{d}p}{\mathrm{d}n}
\lesssim \Rn \cdot c_{\mathrm{BI}} \sqrt{p(n)},
\end{equation}
where $\Rn$ is a scalar redundancy functional satisfying
$\Rn \sim p(n)^{-\alpha}$ in the pre-saturation window.
Separation of variables yields $p(n) \sim n^{1/(\alpha + 1/2)}$,
giving $\bstar = 1/(\alpha + 1/2)$.

The question is: \emph{does the fixed-degree Heisenberg cascade supply a native object
identifiable with~$\Rn$, together with the remaining premises of~\eqref{eq:growth}?}

\subsection{Main results}
\label{sec:mainresults}

\begin{enumerate}[label=(\arabic*)]
\item We audit the O3--O7 derivation chain and identify the three
precise points where the scalar-observable hypothesis enters
(Section~\ref{sec:audit}).
\item We prove a non-equivalence theorem: the block-averaged exponent
$\dexact$ does not coincide with the scalar exponent~$\alpha$ of the
growth equation under generic conditions on block heterogeneity
(Section~\ref{sec:noneq}).
\item We define the effective block-weighted redundancy $\Reff$ and show what aggregation
can and cannot change
(Section~\ref{sec:Reff}).
\item We exhibit the native identity $\Delta r_n=|S_n|\Sigma_n$ and show why the LPS
Born--Infeld growth factor does not transfer (Section~\ref{sec:rederivation}).
\item We retain C1 and C2 only as legacy phenomenological diagnostics and specify the
block-level measurement needed to quantify aggregation (Section~\ref{sec:discrimination}).
\end{enumerate}

\section{Audit of the O3--O7 Derivation Chain}
\label{sec:audit}

We trace the logical chain that produces $\bstar = 1/(\delta + 1/2)$
and identify, at each step, the hypotheses that hold for the proxy
observable but fail for the exact Weil-block observable.

\subsection{Step A: O3 mass formula}
\label{sec:stepA}

O3~\cite{Beau2026a7} derives the inter-generation mass ratio
\begin{equation}
\label{eq:massratio}
\frac{M_i}{M_j}
= \left(\frac{|{\lambda_i - 1}|}{|{\lambda_j - 1}|}\right)^{2/\beta},
\end{equation}
from the support-narrowing mechanism under the power-law growth
hypothesis $p(n) \sim A n^\beta$.
The derivation uses two hypotheses:
\begin{description}[style=nextline]
\item[H-O3-1] The growth $p(n) \sim A n^\beta$ is a power law.
\item[H-O3-2] The mass is inversely proportional to the exit rank:
$M_i \propto E_P / n_{\mathrm{exit}}(\lambda_i)$.
\end{description}
Both hypotheses are explicit structural postulates of O3.
Neither involves any property of the capacity observable.
The mass formula~\eqref{eq:massratio} is therefore \emph{independent}
of the observable-class mismatch and remains valid throughout.

\subsection{Step B: O6 growth equation}
\label{sec:stepB}

O6~\cite{Beau2026a10}, Proposition~4.2, derives $\bstar = 1/(\alpha + 1/2)$
from the growth equation~\eqref{eq:growth}.
The derivation involves three implicit hypotheses about the
redundancy functional~$\Rn$:

\begin{description}[style=nextline]
\item[H-O6-1 (Scalar observable)]
$\Rn$ is a single real-valued function of the cascade rank~$n$,
independent of any internal representation label.
\item[H-O6-2 (Power-law decay)]
$\Rn \sim p(n)^{-\alpha}$ for a single exponent~$\alpha$ over the
pre-saturation window.
\item[H-O6-3 (Direct coupling)]
$\Rn$ enters the growth equation multiplicatively:
the rate of valence increase is proportional to~$\Rn$.
\end{description}

Hypothesis H-O6-1 is the central point of failure.
In the proxy setting (O7--O11), the capacity $C_n^{\mathrm{proxy}}$
is defined via a Fourier-character proxy that abelianises the
representation structure~\cite{Beau2026a11}.
It is indeed a scalar function of~$n$.

In the exact setting (O12--O14), the observable is
$\Sn = \mathbb{E}_c[\Sigma_n^{(c)}]$, a mean over $(q-1)$ Weil blocks,
each of internal dimension~$q$.
This is not a scalar observable in the sense of H-O6-1: it carries
a hidden dependence on the block structure through the averaging.

\subsection{Step C: O7 reinterpretation}
\label{sec:stepC}

O7~\cite{Beau2026a11} recasts the O6 formula in capacity language.
The identification $\alpha = \delta$ (O7, Section~5.1) holds in the
reduced filling model where $p(n) = n$.
The state law $R_n^{(k)} \approx \Phi(\eta_n)$ with
$\Phi(\eta) = 1/\sqrt{1 + \eta^2}$ is derived exactly in the reduced
model and is a property of the scalar capacity.

O7 does not introduce new deductive content for the $\delta \to \bstar$
relation.
It inherits H-O6-1 from O6 without modification.
The Remark~5.2 of O7 states this explicitly:
the formula ``does not introduce a new derivation of~$\beta$:
it recasts the O6 target in capacity language without circularity.''

\subsection{Summary of audit}
\label{sec:auditsummary}

\begin{table}[ht]
\centering
\caption{Audit of the derivation chain.
The column ``Exact status'' indicates whether the hypothesis
holds for the exact Weil-block observable.}
\label{tab:audit}
\begin{tabular}{llcc}
\toprule
Paper & Hypothesis & Proxy & Exact \\
\midrule
O3 & H-O3-1 (power-law growth) & valid & valid \\
O3 & H-O3-2 (mass $\propto 1/n_{\mathrm{exit}}$) & valid & valid \\
O6 & H-O6-1 (scalar observable) & valid & \textbf{fails} \\
O6 & H-O6-2 (power-law decay) & valid & valid \\
O6 & H-O6-3 (direct coupling) & valid & open \\
O7 & inherits H-O6-1 & valid & \textbf{fails} \\
\bottomrule
\end{tabular}
\end{table}

The mass formula of O3 is intact.
The growth equation of O6 is algebraically correct in its domain.
The failure is localised at H-O6-1: the scalar-observable hypothesis,
which enters the growth equation through the identification of~$\Rn$
with a single power-law observable.

\section{Non-Equivalence of Block-Averaged and Scalar Exponents}
\label{sec:noneq}

We now prove that the exponent extracted from a block-averaged observable
does not, in general, coincide with the exponent of a scalar observable
satisfying the growth equation.

\subsection{Setup}

Fix an odd prime~$q$.
The Heisenberg group $G_q = \mathrm{Heis}_3(\mathbb{Z}/q\mathbb{Z})$
has $(q-1)$ faithful irreducible representations
$\rho_c$ ($c \in (\mathbb{Z}/q\mathbb{Z})^\times$), each of dimension~$q$.
The block-level projective capacity at BFS depth~$n$ in block~$c$ is
\begin{equation}
\label{eq:blockcap}
\Sigma_n^{(c)} = \frac{\Delta r_n^{(c)}}{|S_n|},
\end{equation}
where $\Delta r_n^{(c)}$ is the incremental projective span in
representation~$\rho_c$ and $|S_n|$ is the BFS shell size.
The block average is $\Sn = \mathbb{E}_c[\Sigma_n^{(c)}]$.

In the proxy setting, the capacity is a single scalar
$C_n^{\mathrm{proxy}}$ defined via the Fourier-character
abelianisation~\cite{Beau2026a15}.

\subsection{Block heterogeneity and exponent bias}

\begin{definition}[Block exponent and mean exponent]
\label{def:exponents}
For each block~$c$, define the block exponent
$\delta^{(c)}$ by the power-law fit
$\Sigma_n^{(c)} \sim A_c \, n^{-\delta^{(c)}}$
over the fitting window $[n_0, n_1]$.
Define the mean exponent
$\bar{\delta} = \mathbb{E}_c[\delta^{(c)}]$
and the exponent of the mean
$\dexact$ by the fit
$\Sn \sim A \, n^{-\dexact}$.
\end{definition}

\begin{proposition}[Non-equivalence]
\label{prop:noneq}
If the block amplitudes $A_c$ are not all equal,
then $\dexact \neq \bar{\delta}$ in general.
Specifically, let $V_A = \mathrm{Var}_c(A_c) / \mathbb{E}_c[A_c]^2$
be the relative amplitude variance.
Then
\begin{equation}
\label{eq:noneq}
\dexact = \bar{\delta}
- \frac{1}{\log(n_1/n_0)}
\log\!\left(
\frac{\mathbb{E}_c[A_c \, (n_1/n_0)^{-\delta^{(c)}}]}
{\mathbb{E}_c[A_c] \cdot (n_1/n_0)^{-\bar{\delta}}}
\right).
\end{equation}
The correction term vanishes if and only if $\delta^{(c)}$ is independent
of~$c$ (block universality).
\end{proposition}

\begin{proof}
Write $\Sn = \mathbb{E}_c[A_c \, n^{-\delta^{(c)}}]$.
At two endpoints $n_0, n_1$ of the fitting window:
\[
\dexact
= -\frac{\log(\Sn[n_1] / \Sn[n_0])}{\log(n_1 / n_0)}
= -\frac{1}{\log(n_1/n_0)}
\log\!\left(
\frac{\mathbb{E}_c[A_c \, n_1^{-\delta^{(c)}}]}
{\mathbb{E}_c[A_c \, n_0^{-\delta^{(c)}}]}
\right).
\]
If all block exponents equal $\delta_0$, this reduces to
$\delta_0=\bar\delta$.
Otherwise, the ratio is a weighted average of
\[
  \exp\!\left(-\delta^{(c)}\log(n_1/n_0)\right)
\]
with weights $A_c$.
By Jensen's inequality it differs from the value at $\bar\delta$ whenever the distribution of
$\delta^{(c)}$ is non-degenerate.
\end{proof}

\begin{remark}[Direction of the bias]
\label{rem:bias}
For a convex function $f(\delta) = \exp(-\delta \log(n_1/n_0))$,
Jensen's inequality gives
$\mathbb{E}_c[f(\delta^{(c)})] \geq f(\bar{\delta})$,
so $\dexact \leq \bar{\delta}$.
The block-averaged exponent is systematically \emph{smaller}
than the mean of individual block exponents.
This is consistent with the O12--O13 observation that $\dexact$
lies below the target window.
\end{remark}

\begin{remark}[Quantitative estimate from O12--O13 data]
\label{rem:quantitative}
O12~\cite{Beau2026a16} reports inter-block variance ratios
$V_n^{\max}(q)$ that decrease from $\sim 0.15$ at $q = 29$ to
$\sim 0.03$ at $q = 211$.
Even at $q = 211$, the residual variance produces a bias of order
$|\dexact - \bar{\delta}| \sim 0.1$--$0.3$, which is insufficient
to bridge the gap to the target $\delta \in [7.4, 10.6]$.
The non-equivalence is therefore real but quantitatively small
compared to the total deficit.
\end{remark}

\subsection{Consequence for the growth equation}

The growth equation~\eqref{eq:growth} couples the rate of valence
increase to the redundancy of new relational directions.
Physically, it is the \emph{fraction of genuinely new directions
contributed by the BFS frontier} that limits the growth of~$p$.
This fraction is a property of the \emph{aggregate} projective
dynamics, not of any individual block.

The scalar exponent~$\alpha$ in H-O6-2 is the exponent of this
aggregate fraction.
It is not, in general, equal to $\dexact$ (the exponent of the
block-averaged capacity), because the block-averaged capacity
includes internal normalisation and phase structure that do not
enter the growth equation.

This motivates the definition of the effective dynamical redundancy
in the next section.

\section{Effective Dynamical Redundancy}
\label{sec:Reff}

\subsection{Physical meaning of the growth equation}

The growth equation~\eqref{eq:growth} has a precise physical content:
at each cascade step~$n$, the substrate attempts to increase its
effective relational valence~$p$ by incorporating new relational
directions from the BFS frontier.
The rate of successful incorporation is limited by two factors:
\begin{enumerate}[label=(\alph*)]
\item The Born--Infeld flux bound $c_{\mathrm{BI}}$, which caps the rate
of change of any relational variable~\cite{Beau2026a8}.
\item The fraction of frontier directions that are genuinely
new, i.e.\ not already represented in the current relational
structure.
\end{enumerate}
Factor~(b) is what $\Rn$ encodes.
In the proxy setting, where the capacity is a scalar, $\Rn$ is
directly proportional to the capacity: $\Rn \propto \Sigma_n$.
In the exact setting, the identification is more subtle because
``genuinely new'' must be defined relative to the full block structure.

\subsection{Definition of $\Reff$}

\begin{definition}[Effective dynamical redundancy]
\label{def:Reff}
The effective dynamical redundancy at cascade rank~$n$ is
\begin{equation}
\label{eq:Reffdef}
\Reff = \frac{1}{q-1} \sum_{c \in (\mathbb{Z}/q\mathbb{Z})^\times}
w_c(n) \, \Sigma_n^{(c)},
\end{equation}
where $w_c(n)$ is the dynamical weight of block~$c$, defined as
\begin{equation}
\label{eq:weights}
w_c(n)
= \frac{\Delta r_n^{(c)} / \dim V_c}
{\sum_{c'} \Delta r_n^{(c')} / \dim V_{c'}}
= \frac{\Delta r_n^{(c)} / q}
{\sum_{c'} \Delta r_n^{(c')} / q}
= \frac{\Delta r_n^{(c)}}{\sum_{c'} \Delta r_n^{(c')}}.
\end{equation}
\end{definition}

The weight $w_c(n)$ measures the relative contribution of block~$c$
to the total incremental span at step~$n$.
It is the natural dynamical weight because the growth equation
counts the total number of new relational degrees of freedom,
and each block contributes $\Delta r_n^{(c)}$ such degrees.

\begin{proposition}[Relation to the unweighted mean]
\label{prop:relation}
Let $\Sn = \mathbb{E}_c[\Sigma_n^{(c)}]$ be the unweighted block
average (the O12/O13 observable) and let
$\Sigma_n^{w} = \sum_c w_c \Sigma_n^{(c)}$ be the weighted observable
of Definition~\ref{def:Reff}.
Then
\begin{equation}
\label{eq:relation}
\Reff = \Sn + \mathrm{Cov}_c\!\left(
\frac{w_c}{1/(q-1)}, \, \Sigma_n^{(c)}
\right),
\end{equation}
where the covariance is taken over the uniform distribution on blocks.
The correction vanishes if the weights are uniform ($w_c = 1/(q-1)$
for all~$c$), which occurs if and only if all blocks contribute
equally to the incremental span.
\end{proposition}

\begin{proof}
Direct computation:
$\Reff - \Sn
= \sum_c w_c \Sigma_n^{(c)} - \frac{1}{q-1} \sum_c \Sigma_n^{(c)}
= \sum_c (w_c - \frac{1}{q-1}) \Sigma_n^{(c)}
= (q-1) \mathrm{Cov}_c(w_c, \Sigma_n^{(c)})$.
\end{proof}

\begin{remark}[Sign of the correction]
\label{rem:sign}
Blocks with larger $\Delta r_n^{(c)}$ (hence larger $w_c$) tend to
have larger $\Sigma_n^{(c)}$ (since $\Sigma_n^{(c)}$ is the ratio
$\Delta r_n^{(c)} / |S_n|$).
The covariance is therefore generically positive, giving
$\Reff \geq \Sn$.
The effective dynamical redundancy is \emph{larger} than the
unweighted block average.
\end{remark}

\subsection{Exponent of $\Reff$}

\begin{definition}[Block-weighted aggregation exponent]
\label{def:alphaeff}
The block-weighted aggregation exponent~$\alpha_{\mathrm{dyn}}$ is defined by the
power-law fit $\Reff \sim B \, n^{-\alpha_{\mathrm{dyn}}}$ over the
pre-saturation window.
\end{definition}

By Proposition~\ref{prop:relation} and Remark~\ref{rem:sign},
$\alpha_{\mathrm{dyn}} \leq \dexact$: the aggregation exponent is
smaller than or equal to the unweighted block-average exponent.

\begin{corollary}[Aggregation no-go]
\label{cor:nogo}
No choice of dynamical weights $w_c(n)$ satisfying $w_c \geq 0$
and $\sum_c w_c = 1$ can produce
$\alpha_{\mathrm{dyn}} > \dexact$.
In particular, the block-averaged exponent $\dexact \in [3.6, 4.8]$
measured in O12--O13 is an \emph{upper bound} on the decay exponent
of the block-weighted aggregation observable.
Since $\dexact < 5.0 \ll 7.4$, no aggregation or reweighting
of the existing block-level data can bring that aggregation
exponent into the target window $[7.4, 10.6]$.
\end{corollary}

\begin{proof}
At fixed BFS depth~$n$, the shell size $|S_n|$ is common to all blocks.
Therefore $\Sigma_n^{(c)} = \Delta r_n^{(c)} / |S_n|$ is a monotone
increasing function of $\Delta r_n^{(c)}$ alone.
The weight $w_c = \Delta r_n^{(c)} / \sum_{c'} \Delta r_n^{(c')}$
is likewise a monotone increasing function of the same variable
$\Delta r_n^{(c)}$.
Since both $w_c$ and $\Sigma_n^{(c)}$ are monotone increasing functions
of a single common variable, the Chebyshev sum inequality gives
\[
\sum_c w_c \, \Sigma_n^{(c)}
\;\geq\;
\left(\sum_c w_c\right)
\left(\frac{1}{q-1} \sum_c \Sigma_n^{(c)}\right)
= \Sn,
\]
so $\Reff \geq \Sn$ pointwise in~$n$.
A pointwise lower bound on the observable translates into an upper
bound on the decay exponent: $\alpha_{\mathrm{dyn}} \leq \dexact$.
\end{proof}

This is the sharpest negative result of the paper:
it eliminates the entire class of solutions based on
``better averaging of the blocks.''
No refined treatment of the block statistics supplies a native capacity-to-rate bridge.
Any future resolution requires a new growth carrier or an explicitly stated cross-substrate
hypothesis.

\section{Audit of the Transplanted Growth Equation}
\label{sec:rederivation}

\subsection{Native identity and the two-substrate ansatz}

In the exact setting, the BFS frontier at step~$n$ contributes new
relational degrees of freedom through all $(q-1)$ Weil blocks
simultaneously.
The total incremental span is
\begin{equation}
\label{eq:totalspan}
\Delta r_n^{\mathrm{tot}}
= \sum_{c} \Delta r_n^{(c)}.
\end{equation}
Each block~$c$ lives in a representation space of dimension~$q$,
so the total available representation space has dimension
$(q-1) \cdot q \sim q^2$.

Importing the LPS Born--Infeld growth factor would give the two-substrate ansatz
\begin{equation}
\label{eq:growthexact}
\frac{\mathrm{d}p}{\mathrm{d}n}
\lesssim c_{\mathrm{BI}} \sqrt{p(n)} \cdot
\frac{\Delta r_n^{\mathrm{tot}}}{(q-1) \cdot q},
\end{equation}
where the denominator $(q-1) \cdot q$ normalises by the total number
of representation degrees of freedom.
This equation is not native to the Heisenberg cascade: its
$c_{\mathrm{BI}}\sqrt{p(n)}$ factor comes from the changing-degree expander model, while
the exact-block factor comes from the fixed-degree nilpotent model.
No corpus result couples the two~\cite{Beau2026sgn}.

\begin{remark}[Status of the $q$-dependent normalisation]
\label{rem:qnorm}
The factor $1/q$ in~\eqref{eq:growthexact} is a
\emph{dimensional ansatz}, not a derived result.
Its motivation is as follows: in the proxy setting, the capacity
is defined in a scalar (one-dimensional) representation, so no
internal dimension enters the normalisation.
In the exact setting, each block has $q$ internal degrees of
freedom, and the fraction of genuinely new directions per
available degree of freedom naturally acquires a factor $1/q$.
O14 used this factor in an endpoint diagnostic, but proved on revision that a constant
$q^{-\eta}$ cannot change an intra-$q$ regression slope.
However, the Born--Infeld constraint $|\partial_t \chi_v| \leq c_{\mathrm{BI}}$
operates at the level of individual relational variables, not at
the representation level.
The correct translation of this per-variable bound into a
per-representation constraint depends on how the representation
degrees of freedom couple to the relational dynamics,
which remains an open derivation.
The factor $1/q$ is the simplest dimensional estimate for comparing amplitudes across blocks;
a specified inter-$q$ estimator might instead use $q^{-\eta}$.
Neither choice supplies the missing Born--Infeld coupling.
\end{remark}

\subsection{Separation of variables}

Rewrite~\eqref{eq:growthexact} using
$\Delta r_n^{\mathrm{tot}} / ((q-1) \cdot q)
= \frac{1}{q} \cdot \frac{1}{q-1} \sum_c \Delta r_n^{(c)} / |S_n|
\cdot |S_n|
= \frac{|S_n|}{q} \cdot \Sn$,
where we used $\Sn = \frac{1}{q-1} \sum_c \Sigma_n^{(c)}$
with $\Sigma_n^{(c)} = \Delta r_n^{(c)} / |S_n|$.
Thus the exact Heisenberg bookkeeping identity is
\begin{equation}
\label{eq:growthexact2}
\frac{1}{q-1}\Delta r_n^{\mathrm{tot}} = |S_n|\Sn.
\end{equation}

In the BFS growth regime on the Heisenberg graph,
$|S_n| \sim C_S \, n^{D-1}$ with $D$ the effective growth dimension,
and $\Sn \sim A \, n^{-\dexact}$.
The factor $|S_n| \cdot \Sn \sim C \, n^{D - 1 - \dexact}$.
The shell term $(D-1)$ is additive in this increment law.

For the proxy derivation (O6), $|S_n|$ is absorbed into the
definition of~$\Rn$ and does not appear explicitly.
The exact identity makes the shell growth and the capacity decay visible as separate factors,
but contains neither $p(n)$ nor $c_{\mathrm{BI}}$.

\begin{proposition}[Failure of the legacy reciprocal re-derivation]
\label{hyp:growthexact}
The exact identity~\eqref{eq:growthexact2} does not imply a law for $p(n)$.
Even if the two-substrate ansatz~\eqref{eq:growthexact} is imposed and treated as saturated,
direct matching of powers does not give the legacy expression
\begin{equation}
\label{eq:betaexact}
\beta \stackrel{\mathrm{legacy}}{=}
\frac{1}{\tfrac{1}{2} + \dexact - (D - 1)
+ \log q / \log n^{*}},
\end{equation}
because the fixed factor $q^{-1}$ changes no power of $n$.
The saturated exponent balance would instead give
$\beta=2(D-\dexact)$ under the imported square-root factor.
Without saturation of the inequality, even that equality does not follow.
\end{proposition}

\begin{remark}[Native law and scope]
\label{rem:status}
If $\Sn\asymp n^{-\dexact}$ and $|S_n|\asymp n^{D-1}$ on a finite window, the native identity
gives $\Delta r_n\asymp n^{D-1-\dexact}$ and the corresponding continuous cumulative scaling
$r_n\asymp n^{D-\dexact}$ on that window.
This effective law is subject to the integer-rank and saturation caveats established in the
Span-Growth Note~\cite{Beau2026sgn}.
It is a law for span growth, not for the LPS valence $p(n)$ or the phenomenological $\bstar$.
\end{remark}

\begin{proof}
Substitute $p(n)=A n^\beta$, $|S_n|\Sn\asymp n^{D-1-\dexact}$ into the additional
ansatz~\eqref{eq:growthexact}, and compare powers after treating the inequality as saturated:
\[
A \beta n^{\beta - 1}
\asymp \frac{C}{q}c_{\mathrm{BI}}\sqrt{A}\,
n^{\beta/2+D-1-\dexact}.
\]
Thus $\beta-1=\beta/2+D-1-\dexact$, so $\beta=2(D-\dexact)$.
The constant $q^{-1}$ enters only the amplitude.
Equation~\eqref{eq:betaexact} therefore does not follow.
\end{proof}

\subsection{Two legacy phenomenological diagnostics}
\label{sec:scenarios}

\begin{description}[style=nextline]
\item[Diagnostic C1: imported reciprocal form]
Historically, the shell growth exponent $D-1$ and the endpoint term
$\log q/\log n^{*}$ were combined into
$\deff = \dexact - (D - 1) + \log q / \log n^{*}$,
and inserted into
\begin{equation}
\label{eq:C1}
\bstar = \frac{1}{\deff + \tfrac{1}{2}},
\end{equation}
This is a phenomenological diagnostic, not a consequence of the exact growth identity.
The $(D-1)$ term belongs additively to the native increment law, while the fixed-$q$
normalisation changes only an intercept.

For the Heisenberg graph with $D \approx 4$ (Bass--Guivarc'h
growth exponent), $(D-1) = 3$, and
$\log q / \log n^{*} \approx 1$ at $n^{*} \sim q$,
one obtains $\deff \approx \dexact - 3 + 1 = \dexact - 2$.
With $\dexact \in [3.6, 4.8]$, this gives
$\deff \in [1.6, 2.8]$, yielding $\bstar \in [0.26, 0.48]$.
This is still above the target $(0.09, 0.13)$.

\item[Diagnostic C2: additional denominator term]
The second historical diagnostic appends a $q$-dependent term:
\begin{equation}
\label{eq:C2}
\bstar = \frac{1}{\deff + \tfrac{1}{2} + \sigma(q)},
\end{equation}
where $\sigma(q)$ quantifies the displacement required to reach a chosen target.
No Heisenberg growth process carrying this term is derived.

For C2 to produce $\bstar \in (0.09, 0.13)$ with
$\deff \in [1.6, 2.8]$, the correction must satisfy
$\sigma(q) \in [5.1, 8.6]$, which would require
$\sigma(q) \sim q^{\kappa}$ for some $\kappa > 0$.
The numerical value is therefore a target diagnostic, not evidence for a
representation-theoretic correction.
\end{description}

\section{Aggregation Measurement}
\label{sec:discrimination}

\subsection{What the block measurement can establish}

A measurement of $\Reff$ (Definition~\ref{def:Reff}) at multiple primes can quantify the
aggregation bias and test its stability with $q$.
It cannot discriminate between native capacity-to-rate laws, because neither C1 nor C2 is
derived as a Heisenberg growth equation.

\subsection{Required measurement: block-level redundancy}

The O12/O13 pipeline~\cite{Beau2026a16,Beau2026a17} computes
$\Sigma_n^{(c)}$ per block but discards the complex fingerprint
vectors $z_b^{(c)}(n)$ after Gram--Schmidt orthogonalisation.
The weights $w_c(n)$ of Definition~\ref{def:Reff} are computable
from the existing pipeline outputs ($\Delta r_n^{(c)}$ is stored).

The block-level redundancy measurement therefore requires:
\begin{enumerate}[label=(\roman*)]
\item Compute $w_c(n)$ from the stored $\Delta r_n^{(c)}$ values.
\item Form $\Reff = \sum_c w_c(n) \Sigma_n^{(c)}$.
\item Fit $\Reff \sim B \, n^{-\alpha_{\mathrm{dyn}}}$ and extract
$\alpha_{\mathrm{dyn}}$.
\item Compare $\alpha_{\mathrm{dyn}}$ with $\dexact$ to measure the
aggregation bias.
\item Compare the result with the reciprocal diagnostics only as a phenomenological check.
\end{enumerate}

Step~(i) is immediately feasible from existing O12/O13 data.
Steps~(ii)--(v) constitute the minimal computational programme for resolving the aggregation
observable; they do not restore the transfer.

\section{Asymptotic Analysis}
\label{sec:asymptotic}

\subsection{The $q \to \infty$ limit}

In the limit $q \to \infty$, three simplifications occur:
\begin{enumerate}[label=(\alph*)]
\item Block universality: $\mathrm{Var}_c(\Sigma_n^{(c)}) \to 0$
by the O13 variance reduction result~\cite{Beau2026a17},
so $\Reff \to \Sn$ and $\alpha_{\mathrm{dyn}} \to \dexact$.
\item The legacy endpoint ratio $\log q / \log n^{*}$ tends to~1 if $n^{*}\sim q$.
\item The O14 variance ansatz models $\epsilon(q)$ as vanishing under bounded variance.
\end{enumerate}
In this limit, the legacy diagnostics reduce to
$\beta_{\mathrm{C1}}^{\mathrm{diag}}=1/(1/2+\dexact-2)$ and
$\beta_{\mathrm{C2}}^{\mathrm{diag}}
=1/(1/2+\dexact-2+\sigma(\infty))$.

These are algebraic comparisons only; no native conclusion follows from either limit.

\subsection{Finite-$q$ phenomenology}

At accessible primes ($q \leq 211$), the correction terms are
non-negligible.
Table~\ref{tab:scenarios} reports the reciprocal diagnostics for the five O12/O13 primes.

\begin{table}[ht]
\centering
\caption{Legacy reciprocal diagnostics C1 and C2 for the five O12/O13 primes.
C1 uses $\deff = \dexact - 2$;
C2 requires $\sigma(q)$ to reach the target.}
\label{tab:scenarios}
\begin{tabular}{ccccc}
\toprule
$q$ & $\dexact$ & $\deff$ (C1) & $\beta_{\mathrm{C1}}^{\mathrm{diag}}$
& $\sigma(q)$ needed (C2) \\
\midrule
29 & 4.37 & 2.37 & 0.348 & 5.4 \\
61 & 4.80 & 2.80 & 0.303 & 5.9 \\
101 & 4.28 & 2.28 & 0.360 & 5.3 \\
151 & 3.90 & 1.90 & 0.417 & 5.0 \\
211 & 3.59 & 1.59 & 0.478 & 4.7 \\
\bottomrule
\end{tabular}
\end{table}

Under C1, the reciprocal diagnostic lies between $0.30$ and $0.48$ for all tested primes.
It is outside the interval $(0.09,0.13)$.

Under C2, the denominator displacement required to hit the target is
$\sigma(q)\in[4.7,5.9]$.
This number has no derived structural interpretation.

\section{Open Directions}
\label{sec:open}

\begin{description}[style=nextline]
\item[O15-O1 (Block-level redundancy extraction)]
Compute $\Reff$ from the O12/O13 pipeline outputs and extract
$\alpha_{\mathrm{dyn}}$ at the five primes
\[
  q \in \{29,61,101,151,211\}.
\]
This is the minimal measurement required to quantify the aggregation bias.

\item[O15-O2 (Shell growth exponent)]
The effective growth dimension $D \approx 4$ (Bass--Guivarc'h)
is an asymptotic estimate.
The effective value of $D$ in the pre-saturation window
$[n_0, n_1]$ at finite~$q$ may differ.
Extracting $D$ directly from the BFS data would sharpen the native increment law.

\item[O15-O3 (New growth carrier)]
Construct a native Heisenberg growth process carrying the pair observable, or state a new
cross-substrate hypothesis with all coupling premises explicit.

\item[O15-O4 (Quark and neutrino sectors)]
Any reciprocal diagnostic in the quark and neutrino sectors must remain labelled as
phenomenological unless a native carrier is constructed.
\end{description}

\section{Conclusion}
\label{sec:conclusion}

The O3--O7 derivation chain that produces
$\bstar = 1/(\delta + 1/2)$ is conditional in its domain:
a scalar global capacity on proxy-level LPS dynamics.
The failure is not an error but a scope limitation:
the scalar-observable hypothesis H-O6-1 does not hold for the
exact Weil-block capacity measured in O12--O14.

The non-equivalence theorem (Proposition~\ref{prop:noneq})
establishes that the block-averaged exponent $\dexact$ does not
coincide with the scalar exponent~$\alpha$ entering the growth
equation, and that the bias goes in the wrong direction
(Remark~\ref{rem:bias}).

The effective block-weighted redundancy $\Reff$
(Definition~\ref{def:Reff}) is a well-defined aggregation observable.
Its exponent $\alpha_{\mathrm{dyn}}$ is smaller than $\dexact$,
so the aggregation correction alone cannot resolve the tension.

The aggregation no-go (Corollary~\ref{cor:nogo}) is the sharpest
negative result: no reweighting of the block-level data can raise
the aggregation exponent into the target window.
No aggregation rule supplies the missing growth carrier.

The exact native relation is $\Delta r_n=|S_n|\Sigma_n$.
It places the shell exponent $(D-1)$ additively in the increment law.
The block factor $q^{-\eta}$ changes no intra-$q$ slope.
Proposition~\ref{hyp:growthexact} further shows that even the explicitly transplanted
square-root equation does not yield the legacy reciprocal expression.
Thus C1 and C2 are phenomenological diagnostics, not alternative native laws.

The minimal next step is the block-level redundancy extraction
(O15-O1): computing $\Reff$ from existing O12/O13 data to measure
$\alpha_{\mathrm{dyn}}$ directly and quantify the aggregation bias.
This measurement is feasible with the current pipeline outputs and
constitutes the critical path forward.

The O-series has traced the chain from obstruction identification
(O8--O9) through exact measurement (O10--O13), theoretical
observable separation (O14), and derivation audit (O15).
The status of Source~S2 is transformed:
it is no longer an unexplained numerical tension,
nor a theoretical impasse,
but a precise no-transfer result plus a separate aggregation measurement.
The missing object is a native pair-capacity growth carrier, not a reweighting of the existing
block data~\cite{Beau2026sgn}.

\bibliographystyle{plain}
      \bibliography{cosmochrony-bibliography}

\end{document}
